\documentclass[11pt,a4paper]{article}
\usepackage[margin=25mm,headheight=15pt]{geometry}
\usepackage[T1]{fontenc}
\usepackage[utf8]{inputenc}
\IfFileExists{libertinus.sty}{\usepackage{libertinus}}{\usepackage{lmodern}}
\usepackage{microtype,amsmath,amssymb,amsthm,mathtools,mathrsfs}
\usepackage{booktabs,longtable,array,tabularx}
\usepackage{xcolor,fancyhdr,enumitem,listings,xurl,needspace}
\usepackage{hyperref}
\usepackage[nameinlink,noabbrev]{cleveref}
\definecolor{linkblue}{RGB}{25,65,110}
\definecolor{lightgray}{RGB}{246,247,249}
\hypersetup{colorlinks=true,linkcolor=linkblue,citecolor=linkblue,urlcolor=linkblue,
  pdftitle={Recurrence-Free Dyadic Values of the Fabius and Rvachev Functions},
  pdfauthor={Research article prepared for Vladimir Reshetnikov},
  pdfsubject={Exact finite sums, binary compression, and quarter-base spline extraction}}
\pagestyle{fancy}\fancyhf{}
\fancyhead[L]{\small Recurrence-free Fabius--Rvachev values}
\fancyhead[R]{\small\thepage}
\renewcommand{\headrulewidth}{0.35pt}
\setlength{\emergencystretch}{2.5em}
\setlength{\parskip}{0.2em}
\setcounter{tocdepth}{2}
\allowdisplaybreaks[2]
\numberwithin{equation}{section}
\newtheorem{theorem}{Theorem}[section]
\newtheorem{lemma}[theorem]{Lemma}
\newtheorem{proposition}[theorem]{Proposition}
\newtheorem{corollary}[theorem]{Corollary}
\theoremstyle{definition}
\newtheorem{definition}[theorem]{Definition}
\newtheorem{example}[theorem]{Example}
\theoremstyle{remark}
\newtheorem{remark}[theorem]{Remark}
\newcommand{\up}{\operatorname{up}}
\newcommand{\E}{\mathbb E}
\newcommand{\Pp}{\mathbb P}
\newcommand{\R}{\mathbb R}
\newcommand{\N}{\mathbb N}
\newcommand{\Q}{\mathbb Q}
\newcommand{\eps}{\varepsilon}
\newcommand{\dd}{\,\mathrm d}
\newcommand{\coeff}[2]{[#1^{#2}]}
\newcommand{\floor}[1]{\left\lfloor #1\right\rfloor}
\newcommand{\pos}[1]{\left(#1\right)_+}
\newcommand{\qbinom}[3]{\begin{bmatrix}#1\\#2\end{bmatrix}_{#3}}
\newcommand{\sinc}{\operatorname{sinc}}
\newcommand{\sgn}{\operatorname{sgn}}
\newcommand{\bits}{s_2}
\lstset{basicstyle=\ttfamily\footnotesize,columns=fullflexible,keepspaces=true,
  breaklines=true,showstringspaces=false,frame=single,backgroundcolor=\color{lightgray},
  xleftmargin=0.5em,xrightmargin=0.5em,aboveskip=1em,belowskip=1em}
\title{\bfseries Recurrence-Free Dyadic Values of the\\Fabius and Rvachev Functions\\[0.6em]
\large Finite moment elimination, binary compression,\\and exact quarter-base spline extraction}
\author{Research article prepared for Vladimir Reshetnikov}
\date{4 September 2026}
\begin{document}
\maketitle
\begin{abstract}
Exact rationality at dyadic arguments is not, by itself, a closed evaluation formula: a displayed expression may still depend on recursively computed moments or on an infinite limiting process. This article removes both dependencies. Its principal formula expresses $F(m/2^n)$ through only $\lfloor n/2\rfloor+1$ explicitly specified finite Thue--Morse power sums, with rational interpolation weights at the quarter base $q=1/4$. The reason is an exact, not asymptotic, statement: centered finite-convolution values at a fixed dyadic point are polynomials in $4^{-N}$. A local version identifies the entire shrinking spline cell with a finite deconvolution of the dyadic Taylor polynomial. For reduced interior dyadics, the degree is exactly $\lfloor n/2\rfloor$; the resulting number of samples is optimal for an interpolation rule whose weights work on the whole grid.

A second principal formula compresses the arithmetic to one outer summand for each nonzero binary digit of the numerator. Its coefficients are given explicitly by either positive ordered-composition sums or Bernoulli multiplicity sums, and the Bernoulli numbers themselves are supplied by finite sums. Additional formulas recover the required moments from finitely many uniform convolutions, give a certified rounding representation, evaluate the Rvachev bump and the signed global extension, and extend the extraction mechanism to integer-base geometric-uniform laws. All evaluation formulas are finite and recurrence-free. Complete derivations, worked rational examples, and independently cross-checked exact-arithmetic code accompany the article. No global priority claim or Lean formalization of the additional results is asserted.
\end{abstract}
\vfill
\noindent\textbf{Keywords.} Fabius function; Rvachev up-function; dyadic rational; Thue--Morse sequence; finite convolution; Gaussian binomial coefficient; Bernoulli number; exact extrapolation.
\newpage
\begingroup
\fontsize{10}{11}\selectfont
\setlength{\parskip}{0pt}
\tableofcontents
\endgroup
\newpage

\section{The question and the principal answer}
\label{sec:answer}

\subsection{What counts as a closed form here}
We seek the exact values at $m/2^n$, with integer $m$ and nonnegative integer $n$. A formula will be called \emph{recurrence-free} when its right-hand side contains only explicitly bounded finite sums or products of specified numbers, together with elementary arithmetic. A coefficient-extraction symbol will be used only as an abbreviation for a finite sum that is displayed separately. In particular, a sequence of unknown moments defined by a triangular recurrence is not left as an input to any final evaluation formula.

Three functions must be distinguished. Throughout this article, $F$ is the bounded Fabius distribution function, $\up$ is the compactly supported Rvachev bump, and $\Theta$ is the signed global extension. On their principal domains,
\begin{equation}
 \up(x)=F(1-|x|),\qquad F(x)=\up(x-1)\quad(0\le x\le1).
 \label{eq:fold-intro}
\end{equation}
Precise global conventions and their probabilistic construction appear in \cref{sec:setup}.

The repository exposition already contains the moment-based dyadic formula, binary reduction, reciprocal-power composition formulas, and a half-base $q$-binomial layer \cite{repo}. The classical arithmetic and evaluation background is developed by Arias de Reyna and Haugland \cite{arias2018,haugland}. Reshetnikov's explicit-formula question is a direct antecedent \cite{reshetnikov}; Pinelis gives a finite-convolution limit \cite{pinelis}. Here the main additional organization is exact parity-compressed extraction, its local polynomial explanation, and fully substituted finite coefficients. The moment identities and the idea of binary reduction are not claimed as new. The source comparison concerns the cited text, not an exhaustive priority search.

\subsection{A formula using only factorials, powers, and finite sums}
Write
\begin{equation}
 T_r=\frac{r(r+1)}2,\qquad \eps_h=(-1)^{\bits(h)},
 \label{eq:basic-notation}
\end{equation}
where $\bits(h)$ is the sum of the binary digits of $h$. This sign is not defined by a recurrence: for $0\le h<2^L$, its digits are the explicit integers $\lfloor h/2^\ell\rfloor-2\lfloor h/2^{\ell+1}\rfloor$ for $0\le\ell<L$. Define the \emph{explicit integer power sum}
\begin{equation}
 \mathcal S_N(A)=\sum_{h=0}^{A-1}\eps_h(2A-2h-1)^N
 \qquad(N\ge1,\ A\ge0).
 \label{eq:integer-power-sum}
\end{equation}
An empty sum is zero, and an empty product is one.

\begin{theorem}[Quarter-base closed formula]
\label{thm:main-quarter}
Let $n\ge1$, $0\le m\le2^n$, and $d=\lfloor n/2\rfloor$. Put
\begin{equation}
 w_{d,j}=
 \frac{(-1)^{d-j}4^{T_j}}
 {\displaystyle\prod_{r=1}^{j}(4^r-1)
  \displaystyle\prod_{r=1}^{d-j}(4^r-1)}
 \qquad(0\le j\le d).
 \label{eq:weights-integer}
\end{equation}
Then
\begin{equation}
 \boxed{\displaystyle
 F\!\left(\frac m{2^n}\right)=
 \sum_{j=0}^{d}w_{d,j}\,
 \frac{\mathcal S_{n+j}(2^jm)}{2^{T_{n+j}}(n+j)!}.}
 \label{eq:main-quarter}
\end{equation}
More generally, for any integer $R\ge0$,
\begin{equation}
 F\!\left(\frac m{2^n}\right)=
 \sum_{j=0}^{d}w_{d,j}\,
 \frac{\mathcal S_{n+R+j}(2^{R+j}m)}
 {2^{T_{n+R+j}}(n+R+j)!}.
 \label{eq:main-quarter-shift}
\end{equation}
\end{theorem}

The proof is given in \cref{sec:extraction}. It does \emph{not} assume rationality of the target value. Thus \cref{eq:main-quarter} is itself a proof of dyadic rationality, using only rational finite expressions. It contains no moments, Bernoulli numbers, derivative values, auxiliary recurrent relations, or limits.

The first rows of weights are especially simple:
\begin{equation}
\begin{array}{c|l}
 d & (w_{d,0},\ldots,w_{d,d})\\\hline
 0 &(1)\\
 1 &(-1,4)/3\\
 2 &(1,-20,64)/45\\
 3 &(-1,84,-1344,4096)/2835.
\end{array}
\label{eq:weight-rows}
\end{equation}
They sum to one. These are exact interpolation rows, not fitted numerical coefficients.

\subsection{How to choose among the formulas}
The quarter-base formula is the most self-contained: it uses only explicit integer power sums and rational weights. The binary formula in \cref{sec:binary} is usually more economical when the denominator depth is large, because it replaces a long signed prefix sum by at most $\bits(m)$ blocks. The positive composition coefficients in \cref{sec:compositions} avoid Bernoulli numbers altogether. The partition formula in \cref{sec:cumulants} makes the cumulant structure visible. The finite-prefix formula in \cref{sec:moment-extraction} provides another independent elimination of the same coefficients.

Recursion can of course be a very efficient implementation technique. The point here is different: no recurrence is required to \emph{define} the numbers on the right-hand side of an evaluation identity.

\section{Normalization and the probability model}
\label{sec:setup}

\subsection{The bounded distribution and its centered density}
Let $U_1,U_2,\ldots$ be independent random variables, each uniform on $[0,1]$. Define
\begin{equation}
 X=\sum_{j=1}^{\infty}2^{-j}U_j,
 \qquad F(x)=\Pp(X\le x).
 \label{eq:random-X}
\end{equation}
The sum converges for every choice of the coordinates and takes values in $[0,1]$. Let $V_j=2U_j-1$, so that $V_j$ is uniform on $[-1,1]$. Then
\begin{equation}
 Z=2X-1=\sum_{j=1}^{\infty}2^{-j}V_j.
 \label{eq:random-Z}
\end{equation}
We denote its density by $\up$. This is the standard probability construction underlying the functions \cite{arias1982,arias2018}.

\begin{lemma}[Smoothness and normalization]
\label{lem:normalization}
The density $\up$ is even, nonnegative, $C^\infty$, and supported on $[-1,1]$. Moreover,
\begin{align}
 F'(x)&=2\up(2x-1),\label{eq:F-density}\\
 \up(x)&=F(1-|x|),\label{eq:fold}\\
 F(1-x)&=1-F(x),\label{eq:F-reflection}\\
 \up(0)&=1,\qquad 0\le\up(x)\le1.\label{eq:up-bound}
\end{align}
The bounded $F$ is $C^\infty$ on $\R$, equals zero for $x\le0$ and one for $x\ge1$, and is $2$-Lipschitz.
\end{lemma}
\begin{proof}
The law of $Z$ is symmetric because every $V_j$ has a symmetric law. Its characteristic function is
\[
 \E e^{itZ}=\prod_{j\ge1}\frac{\sin(t/2^j)}{t/2^j}.
\]
For any positive integer $L$, bounding the first $L$ factors by $\min(1,2^j/|t|)$ and all remaining factors by one gives decay $O_L((1+|t|)^{-L})$. Consequently the characteristic function multiplied by any fixed power of $t$ is integrable. Fourier inversion and differentiation under its integral give a $C^\infty$ density. Nonnegativity and support follow from the law of the bounded sum. The affine change $Z=2X-1$ gives \cref{eq:F-density}.

Use the independent-copy decomposition $Z\overset d=(V+Z')/2$, where $V$ is uniform on $[-1,1]$ and $Z'$ has the law of $Z$. The density of $V/2$ is one on $[-1/2,1/2]$, so
\begin{equation}
 \up(x)=\Pp(2x-1\le Z'\le2x+1).
 \label{eq:up-head}
\end{equation}
For $x\ge0$, the upper endpoint is at least one; hence the last probability is $\Pp(X'\ge x)=1-F(x)$. Symmetry of $X$ about $1/2$ gives $1-F(x)=F(1-x)$. For $x\le0$, use evenness. These arguments also apply outside the support, using the constant tails of $F$. Thus \cref{eq:fold} and \cref{eq:F-reflection} follow. Formula \cref{eq:up-head} gives $\up(0)=1$ and $\up\le1$. Finally \cref{eq:F-density} bounds $F'$ by two. Smoothness and the constant exterior values imply flatness at the endpoints.
\end{proof}
In particular $F'(x)=2F(2x)$ on $[0,1/2]$. This differential equation must not be extended to all real $x$ with the \emph{bounded} $F$ on the right.

\subsection{The signed extension and finite derivative jets}
Define the locally finite sum
\begin{equation}
 \Theta(x)=\sum_{b=0}^{\infty}\eps_b\up(x-2b-1).
 \label{eq:Theta}
\end{equation}
On $[2b,2b+2]$ only the $b$th term can be nonzero. In particular $\Theta=F$ on $[0,1]$ and $\Theta$ vanishes on $(-\infty,0]$.

\begin{lemma}[Dilation and dyadic termination]
\label{lem:derivatives}
For every $r\ge0$,
\begin{equation}
 \Theta^{(r)}(x)=2^{T_r}\Theta(2^r x).
 \label{eq:Theta-derivative}
\end{equation}
At a dyadic $x=m/2^n\in[0,1]$,
\begin{equation}
 F^{(r)}(x)=0\qquad(r\ge n+1).
 \label{eq:jet-termination}
\end{equation}
For positive odd $m$ one has
\begin{equation}
 \Theta(m)=\eps_{(m-1)/2},\qquad
 \Theta(m/2)=\frac12\eps_{\lfloor m/4\rfloor}.
 \label{eq:integer-half-values}
\end{equation}
At nonnegative even integers, $\Theta$ vanishes.
\end{lemma}
\begin{proof}
Differentiating \cref{eq:up-head} gives
\[
 \up'(x)=2\up(2x+1)-2\up(2x-1).
\]
Insert this into \cref{eq:Theta}. The two terms combine into the single sum $2\Theta(2x)$ because $\eps_{2b}=\eps_b$ and $\eps_{2b+1}=-\eps_b$. Repeated differentiation yields \cref{eq:Theta-derivative}. Integer values follow immediately from the support of $\up$, $\up(0)=1$, and $\up(\pm1)=0$. At an odd half-integer the argument of the relevant translate is $\pm1/2$, where $\up=F(1/2)=1/2$.

If $r\ge n+1$, the argument $2^r m/2^n$ is an even integer, giving zero. Derivatives of $F$ and $\Theta$ agree at interior points of $[0,1]$; at the endpoints the same assertion follows by smoothness and one-sided agreement. This proves \cref{eq:jet-termination}.
\end{proof}
A terminating Taylor jet does not make the function equal to that polynomial in a fixed neighborhood. The exact spline statement proved below concerns a finite convolution and a cell whose width shrinks with its depth.

\subsection{Moment generating functions and independent tails}
Put
\begin{equation}
 H(t)=\frac{\sinh t}{t},\quad H(0)=1,
 \qquad M(t)=\E e^{tZ}=\prod_{j=1}^{\infty}H(t/2^j).
 \label{eq:MGF}
\end{equation}
The product converges locally uniformly, since $H(t/2^j)=1+O_K(4^{-j})$ on each compact set $K$. It defines an entire even function, also directly obtained from the bounded random variable $Z$. Define the finite prefixes
\begin{equation}
 Z_N=\sum_{j=1}^{N}2^{-j}V_j,
 \qquad M_N(t)=\E e^{tZ_N}=\prod_{j=1}^{N}H(t/2^j),
 \label{eq:prefix-Z}
\end{equation}
with $Z_0=0$ and $M_0=1$. The exact factorization is
\begin{equation}
 M(t)=M_N(t)M(2^{-N}t),\qquad M(2t)=H(t)M(t).
 \label{eq:tail-MGF}
\end{equation}
Quotients by $M$ will only be used near zero, or as formal power series with constant coefficient one. No assertion that $1/M$ is entire is needed.

\section{From a finite box to the dyadic master identity}
\label{sec:master}

\subsection{A finite-convolution formula with proof}
\begin{lemma}[Box-simplex inclusion--exclusion]
\label{lem:box}
Let $a_1,\ldots,a_N>0$ and let $U_j$ be independent uniforms on $[0,1]$. For $N\ge1$,
\begin{equation}
 \Pp\left(\sum_{j=1}^N a_jU_j\le x\right)
 =\frac1{N!\prod_{j=1}^N a_j}
 \sum_{S\subseteq\{1,\ldots,N\}}(-1)^{|S|}
 \pos{x-\sum_{j\in S}a_j}^{N}.
 \label{eq:box}
\end{equation}
\end{lemma}
\begin{proof}
After the substitution $y_j=a_jU_j$, the desired probability is the volume of the part of the box $\prod[0,a_j]$ satisfying $\sum y_j\le x$, divided by $\prod a_j$. The unrestricted nonnegative simplex $\sum y_j\le x$ has volume $x_+^N/N!$: this follows by scaling the unit simplex, whose volume is $1/N!$ by successive integration.

Exclude the conditions $y_j>a_j$ by inclusion--exclusion. For a fixed excluded set $S$, translate $y_j$ by $a_j$ for $j\in S$. The remaining simplex has the shifted bound $x-\sum_{j\in S}a_j$. Summing its volume with sign $(-1)^{|S|}$ proves \cref{eq:box}. Boundaries have zero volume, so the weak inequalities cause no ambiguity.
\end{proof}

For $a_j=2^{-j}$, subset sums are $h/2^N$, with $0\le h<2^N$, and the subset sign is $\eps_h$. If $S_N=\sum_{j=1}^N2^{-j}U_j$ and $H_N$ is its CDF, then
\begin{equation}
 H_N(x)=\frac{2^{T_N}}{N!}\sum_{h=0}^{2^N-1}\eps_h\pos{x-h/2^N}^{N}.
 \label{eq:uncentered-spline}
\end{equation}
This is a finite formula for an approximant, not yet an exact finite formula for $F$.

\subsection{Why the tail becomes polynomial at a dyadic point}
The independent-copy decomposition
\[
 X\overset d=S_n+2^{-n}X'
\]
implies $F(x)=\E H_n(x-2^{-n}X')$. At $x=m/2^n$, the active set in \cref{eq:uncentered-spline} becomes particularly simple: $m-h-X'>0$ almost surely for $h<m$, and it is nonpositive for $h\ge m$. The support $0\le X'\le1$ is exactly what makes this possible.

\begin{theorem}[Dyadic master identity]
\label{thm:master}
For $n\ge1$ and $0\le m\le2^n$,
\begin{align}
 F\!\left(\frac m{2^n}\right)
 &=\frac{2^{-n(n-1)/2}}{n!}
   \sum_{h=0}^{m-1}\eps_h\E(m-h-X)^n\label{eq:master-X}\\
 &=\frac{2^{-T_n}}{n!}
   \sum_{h=0}^{m-1}\eps_h\E(2m-2h-1-Z)^n.
 \label{eq:master-Z}
\end{align}
\end{theorem}
\begin{proof}
Substitute $x=m/2^n$ in the preceding conditional expectation. The factor $2^{T_n}$ from \cref{eq:uncentered-spline} and the scaling $2^{-n^2}$ combine to $2^{-n(n-1)/2}$. Removing the inactive terms gives \cref{eq:master-X}. Now use $X=(Z+1)/2$, which supplies another factor $2^{-n}$ and proves \cref{eq:master-Z}. Symmetry permits either sign in front of $Z$.
\end{proof}

Let
\begin{equation}
 a_k=\coeff{t}{2k}M(t)=\frac{\E Z^{2k}}{(2k)!},
 \qquad
 \Phi_r(v)=\sum_{k=0}^{\lfloor r/2\rfloor}
           a_k\frac{v^{r-2k}}{(r-2k)!}.
 \label{eq:a-Phi}
\end{equation}
Equivalently $\Phi_r(v)=\coeff{t}{r}e^{vt}M(t)$. Thus
\begin{equation}
 F\!\left(\frac m{2^n}\right)
 =2^{-T_n}\sum_{h=0}^{m-1}\eps_h\Phi_n(2m-2h-1).
 \label{eq:master-Phi}
\end{equation}
At this stage the moments have merely been isolated. The next section replaces every $a_k$ by a finite expression. Formula \cref{eq:master-Phi} is also valid for $n=0$, $m\in\{0,1\}$, on setting $\Phi_0=1$.

\section{Eliminating every moment}
\label{sec:moment-elimination}

\subsection{Bernoulli numbers without a recurrence}
\label{sec:cumulants}
Define, for every nonnegative integer $s$,
\begin{equation}
 \beta_s=\sum_{j=0}^{s}\frac1{j+1}
         \sum_{v=0}^{j}(-1)^v\binom jv v^s,
 \label{eq:Bernoulli-finite}
\end{equation}
where $0^0=1$. These are the Bernoulli numbers with convention $\beta_1=-1/2$; the different symbol temporarily emphasizes their finite definition.

\begin{lemma}[Finite Bernoulli identity]
\label{lem:Bernoulli}
As formal power series,
\begin{equation}
 \sum_{s\ge0}\beta_s\frac{z^s}{s!}=\frac{z}{e^z-1}.
 \label{eq:Bernoulli-EGF}
\end{equation}
Consequently $\beta_s=B_s$, with $B_1=-1/2$, and $B_{2r+1}=0$ for $r\ge1$.
\end{lemma}
\begin{proof}
The inner finite difference in \cref{eq:Bernoulli-finite} is zero when $j>s$, as is seen by taking the coefficient of $z^s$ in $(1-e^z)^j$. Therefore the exponential generating function can be rearranged coefficientwise:
\[
 \sum_{j\ge0}\frac{(1-e^z)^j}{j+1}
 =\frac{-\log(1-(1-e^z))}{1-e^z}
 =\frac{z}{e^z-1}.
\]
All operations are formal and locally finite in each degree. Finally
$z/(e^z-1)+z/2=(z/2)\coth(z/2)$ is even, which proves the odd-index vanishing.
\end{proof}

For $r\ge1$, put
\begin{equation}
 \lambda_r=\frac{2^{2r-1}B_{2r}}{r(2r)!(4^r-1)}.
 \label{eq:lambda}
\end{equation}
The values of $B_{2r}$ on this right-hand side can always be replaced by \cref{eq:Bernoulli-finite}.

\begin{proposition}[Explicit cumulants]
\label{prop:log-M}
Near zero, and also formally,
\begin{equation}
 \log M(t)=\sum_{r\ge1}\lambda_rt^{2r}.
 \label{eq:log-M}
\end{equation}
The first coefficients are
\begin{equation}
 \lambda_1=\frac1{18},\quad
 \lambda_2=-\frac1{2700},\quad
 \lambda_3=\frac1{178605},\quad
 \lambda_4=-\frac1{9639000},\quad
 \lambda_5=\frac1{478533825}.
 \label{eq:lambda-small}
\end{equation}
\end{proposition}
\begin{proof}
Differentiate $\log H(t)$:
\[
 \frac{\mathrm d}{\mathrm dt}\log H(t)
 =\coth t-\frac1t
 =1+\frac{2}{e^{2t}-1}-\frac1t.
\]
Insert \cref{eq:Bernoulli-EGF}. The constant term cancels because $B_1=-1/2$, and integration, with zero constant term, gives
\[
 \log H(t)=\sum_{r\ge1}\frac{2^{2r-1}B_{2r}}{r(2r)!}t^{2r}.
\]
Summing this identity at $t/2^j$ over $j\ge1$ multiplies the coefficient of $t^{2r}$ by $\sum_{j\ge1}4^{-rj}=1/(4^r-1)$. Local uniform convergence justifies the analytic argument; coefficientwise geometric summation gives the same formal result.
\end{proof}

\begin{theorem}[Multiplicity-profile coefficients]
\label{thm:Bell}
Set $a_0=1$. For $k\ge1$,
\begin{equation}
 \boxed{\displaystyle
 a_k=\sum_{\substack{v_1,\ldots,v_k\ge0\\
                         v_1+2v_2+\cdots+kv_k=k}}
           \prod_{r=1}^{k}\frac{\lambda_r^{v_r}}{v_r!}.}
 \label{eq:a-partitions}
\end{equation}
Equivalently, for $d=\lfloor n/2\rfloor$,
\begin{equation}
 \boxed{\displaystyle
 F\!\left(\frac m{2^n}\right)
 =2^{-T_n}\sum_{h=0}^{m-1}\eps_h
 \sum_{\substack{v_0,v_1,\ldots,v_d\ge0\\
                         v_0+2\sum_{r=1}^d rv_r=n}}
  \frac{(2m-2h-1)^{v_0}}{v_0!}
  \prod_{r=1}^d\frac{\lambda_r^{v_r}}{v_r!}.}
 \label{eq:F-partitions}
\end{equation}
\end{theorem}
\begin{proof}
Exponentiating \cref{eq:log-M} gives $M(t)=\prod_{r\ge1}\exp(\lambda_rt^{2r})$. To extract degree $2k$, only factors $r\le k$ and powers $v_r\le k/r$ can contribute. Expanding these finitely many exponentials to the relevant orders gives \cref{eq:a-partitions}. Multiplication by $e^{vt}$ gives the inner sum in \cref{eq:F-partitions}; insert this into \cref{eq:master-Phi}.
\end{proof}
This is the complete Bell-polynomial transform written out, so no convention for Bell polynomials is required. It is genuinely finite even when the displayed generating functions are infinite.

\subsection{Positive ordered compositions: no Bernoulli numbers}
\label{sec:compositions}
For a positive integer $k$, an ordered composition of $k$ is a tuple $(r_1,\ldots,r_\ell)$ of positive integers with sum $k$. Define its suffix sums
\begin{equation}
 R_i=r_i+r_{i+1}+\cdots+r_\ell.
 \label{eq:suffix}
\end{equation}

\begin{theorem}[Positive composition formula]
\label{thm:compositions}
For $k\ge1$,
\begin{equation}
 \boxed{\displaystyle
 a_k=\sum_{\ell=1}^{k}
       \sum_{\substack{r_1+\cdots+r_\ell=k\\r_i\ge1}}
       \prod_{i=1}^{\ell}
       \frac1{(2r_i+1)!(4^{R_i}-1)}.}
 \label{eq:a-compositions}
\end{equation}
Every summand is positive. Together with $a_0=1$, this supplies every coefficient in \cref{eq:master-Phi} without any special-number sequence.
\end{theorem}
\begin{proof}
Expand the product \cref{eq:MGF} using
\[
 H(t/2^j)=\sum_{r\ge0}\frac{t^{2r}4^{-jr}}{(2r+1)!}.
\]
In a contribution to degree $2k$, let $j_1<\cdots<j_\ell$ be the distinct occupied levels and let $r_i\ge1$ be the contribution at level $j_i$. Fix the ordered composition $(r_1,\ldots,r_\ell)$. Its sum over occupied levels is
\[
 \frac1{\prod_i(2r_i+1)!}
 \sum_{1\le j_1<\cdots<j_\ell}4^{-\sum_i j_i r_i}.
\]
Put $g_1=j_1$ and $g_i=j_i-j_{i-1}$ for $i>1$. Every $g_i$ is a positive integer and $\sum_i j_ir_i=\sum_i g_iR_i$. Hence the last sum factors into the finite product of geometric sums
\[
 \prod_{i=1}^{\ell}\left(\sum_{g_i\ge1}4^{-g_iR_i}\right)
 =\prod_{i=1}^{\ell}\frac1{4^{R_i}-1}.
\]
The original coefficient sum has nonnegative terms, so its rearrangement is justified by monotone convergence. Summing over all ordered compositions proves the result.
\end{proof}

For example,
\begin{align}
 a_1&=\frac1{3!(4-1)}=\frac1{18},\\
 a_2&=\frac1{5!(16-1)}+
       \frac1{3!(16-1)\,3!(4-1)}
      =\frac{19}{16200}.
 \label{eq:a-small}
\end{align}
Thus $\E Z^2=1/9$ and $\E Z^4=19/675$. Multiplying normalized coefficients by $(2k)!$ is important; omitting that factor would change the dyadic values.

The fully substituted polynomial is
\begin{equation}
\begin{split}
 \Phi_n(v)={}&\frac{v^n}{n!}\\
 &+\sum_{k=1}^{\lfloor n/2\rfloor}\frac{v^{n-2k}}{(n-2k)!}
  \sum_{\ell=1}^{k}
  \sum_{\substack{r_1+\cdots+r_\ell=k\\r_i\ge1}}
  \prod_{i=1}^{\ell}\frac1{(2r_i+1)!(4^{R_i}-1)}.
\end{split}
\label{eq:Phi-compositions}
\end{equation}
Equations \cref{eq:master-Phi,eq:Phi-compositions} therefore form a closed expression made entirely from factorials, integer powers, finite sums, and finite products. The suffix sums $R_i$ are explicit sums of the summation variables, not a recurrent sequence.

\subsection{Finite reciprocal coefficients}
We shall also need
\begin{equation}
 \frac1{M(t)}=\sum_{k\ge0}e_kt^{2k},\qquad e_0=1.
 \label{eq:inverse-M}
\end{equation}
Exactly the same expansion proves the recurrence-free formula
\begin{equation}
 e_k=\sum_{\substack{v_1+2v_2+\cdots+kv_k=k\\v_r\ge0}}
        \prod_{r=1}^{k}\frac{(-\lambda_r)^{v_r}}{v_r!}.
 \label{eq:e-partitions}
\end{equation}
The first values are
\begin{equation}
 e_1=-\frac1{18},\qquad
 e_2=\frac{31}{16200},\qquad
 e_3=-\frac{2347}{42865200},\qquad
 e_4=\frac{1904369}{1311675120000}.
 \label{eq:e-small}
\end{equation}
One may alternatively eliminate $\lambda_r$ by the ordinary finite reciprocal expansion
\begin{equation}
 e_k=\sum_{\ell=1}^{k}(-1)^\ell
      \sum_{\substack{r_1+\cdots+r_\ell=k\\r_i\ge1}}
      \prod_{i=1}^{\ell}a_{r_i}\quad(k\ge1),
 \label{eq:e-compositions}
\end{equation}
and then substitute \cref{eq:a-compositions}. This is the degree-$k$ coefficient of the finite geometric expansion of $1/(1+(M-1))$.

\section{A formula with one outer summand per nonzero binary digit}
\label{sec:binary}

The master formula sums over $m$ consecutive integers. Its special signs permit an exact compression into binary blocks. This is useful independently of the spline-extraction method.

Write the ordinary binary expansion
\begin{equation}
 m=\sum_{j=0}^{n}b_j2^j,\qquad b_j\in\{0,1\},\qquad
 J(m)=\{j:b_j=1\}.
 \label{eq:binary-data}
\end{equation}
The top digit $b_n$ can occur only at the endpoint $m=2^n$. For $j\in J(m)$ define explicitly
\begin{equation}
 \begin{split}
 p_j&=2^{j+1}\floor{\frac{m}{2^{j+1}}},\qquad
 r_j=m-2^j\floor{\frac{m}{2^j}},\\
 A_j&=1+2^{1-j}r_j,\qquad
 \eta_j=(-1)^{\bits(\lfloor m/2^{j+1}\rfloor)}.
 \end{split}
 \label{eq:binary-parameters}
\end{equation}
Thus $p_j$ is the part strictly above the $j$th bit, $r_j$ is the part strictly below it, and $m=p_j+2^j+r_j$. These are digit operations, not recursive values of $F$.

\begin{theorem}[Binary-compressed closed formula]
\label{thm:binary}
For $n\ge0$ and $0\le m\le2^n$,
\begin{equation}
 \boxed{\displaystyle
 F\!\left(\frac m{2^n}\right)
 =\sum_{j\in J(m)}\eta_j\,2^{-T_{n-j}}\Phi_{n-j}(A_j).}
 \label{eq:binary-formula}
\end{equation}
Here $\Phi_r$ is given by the finite composition sum \cref{eq:Phi-compositions}. Equivalently, the entire right-hand side can be written as
\begin{equation}
 \sum_{j\in J(m)}\eta_j2^{-T_{n-j}}
 \sum_{\substack{v_0+2v_1+4v_2+\cdots+2Kv_K=n-j\\v_0,\ldots,v_K\ge0}}
 \frac{A_j^{v_0}}{v_0!}
 \prod_{s=1}^{K}\frac{\lambda_s^{v_s}}{v_s!},
 \qquad K=\floor{\frac{n-j}{2}},
 \label{eq:binary-partition}
\end{equation}
where the inner sum and product have $K$ depending on their outer index, and $\lambda_s$ is explicitly given in \cref{eq:lambda}. There are precisely $\bits(m)$ outer summands.
\end{theorem}

\begin{proof}
The interval of integers $[0,m)$ is the disjoint union of the blocks
$[p_j,p_j+2^j)$ for $j\in J(m)$. There is no binary carry when an integer $h<2^j$ is added to $p_j$, so $\eps_{p_j+h}=\eps_{p_j}\eps_h=\eta_j\eps_h$. The finite Thue--Morse product therefore gives
\begin{equation}
 \sum_{h=0}^{m-1}\eps_h z^h
 =\sum_{j\in J(m)}\eta_j z^{p_j}
       \prod_{\ell=0}^{j-1}(1-z^{2^\ell}).
 \label{eq:prefix-binary-product}
\end{equation}
This product identity also follows simply by expanding the product: selecting a subset of its factors selects the nonzero binary digits of $h$.

From \cref{eq:master-Phi}, the desired value is
\[
 2^{-T_n}\coeff{t}{n}
 e^{(2m-1)t}M(t)\sum_{h<m}\eps_h e^{-2ht}.
\]
For the block of length $2^j$, factor its product as
\begin{equation}
 \prod_{\ell=0}^{j-1}(1-e^{-2^{\ell+1}t})
 =2^{T_j}t^j e^{-(2^j-1)t}
       \prod_{\ell=0}^{j-1}H(2^\ell t).
 \label{eq:binary-H-product}
\end{equation}
Indeed, each factor is $2^{\ell+1}t\,e^{-2^\ell t}H(2^\ell t)$. The infinite product for $M$ gives the exact cancellation
\[
 M(t)\prod_{\ell=0}^{j-1}H(2^\ell t)=M(2^jt).
\]
Consequently this block contributes
\[
 \eta_j2^{-T_n+T_j}
 \coeff{t}{n-j} e^{(2m-2p_j-2^j)t}M(2^jt).
\]
Substituting $s=2^jt$ changes the coefficient by $2^{j(n-j)}$ and changes the exponential argument to
$(2m-2p_j-2^j)/2^j=1+2^{1-j}r_j=A_j$.
Finally,
\[
 -T_n+T_j+j(n-j)=-T_{n-j}.
\]
This proves \cref{eq:binary-formula}. Expanding $e^{A_jt}\exp(\sum_{s\ge1}\lambda_st^{2s})$ proves \cref{eq:binary-partition}. Every coefficient uses only finitely many factors.
\end{proof}

For $m=0$ the empty sum gives zero. For $m=2^n$, its single block gives $\Phi_0(1)=1$. For $m=1$, only the lowest bit occurs, and the formula reduces to
\begin{equation}
 F(2^{-n})=2^{-T_n}\Phi_n(1).
 \label{eq:smallest-dyadic}
\end{equation}
No special endpoint conventions beyond these are needed.

\section{The exact polynomial hidden in centered finite splines}
\label{sec:local-splines}

\subsection{The centered finite distribution}
For $N\ge1$ define
\begin{equation}
 S_N=\sum_{j=1}^{N}2^{-j}U_j,
 \qquad Y_N=S_N+2^{-N-1}=\frac{1+Z_N}{2},
 \qquad G_N(x)=\Pp(Y_N\le x).
 \label{eq:centered-prefix}
\end{equation}
The added midpoint is the mean of the omitted tail. The shift is essential: centering makes the correction polynomial even in its small scale. The uncentered finite CDF does not enjoy this parity simplification.

The finite box formula gives
\begin{equation}
 G_N(x)=\frac{2^{T_N}}{N!}
   \sum_{h=0}^{2^N-1}\eps_h
       \pos{x-\frac{h+1/2}{2^N}}^N.
 \label{eq:G-positive-parts}
\end{equation}
At $x=m/2^n$ and $N\ge n\ge1$, let $A=2^{N-n}m$. Exactly the terms $h<A$ are active, and hence
\begin{equation}
 \boxed{\displaystyle
 G_N\!\left(\frac m{2^n}\right)
 =\frac{1}{2^{T_N}N!}
   \sum_{h=0}^{A-1}\eps_h(2A-2h-1)^N
 =\frac{\mathcal S_N(A)}{2^{T_N}N!}.}
 \label{eq:G-integer-sum}
\end{equation}
The possible negative terms in this expression arise only from inclusion--exclusion; $G_N$ itself is a distribution function.

\subsection{A finite expression for every dyadic derivative}
We first make explicit the Taylor data used in the proof. They will disappear completely from the principal closed formula.

\begin{lemma}[Differentiated master formula]
\label{lem:derivative-master}
For $n\ge1$, $0\le m\le2^n$, and $0\le r\le n$,
\begin{equation}
 F^{(r)}\!\left(\frac m{2^n}\right)
 =2^{(n+1)r-T_n}
   \sum_{h=0}^{m-1}\eps_h\Phi_{n-r}(2m-2h-1).
 \label{eq:derivative-master}
\end{equation}
For $r\ge n+1$ the derivative is zero. In particular, every derivative at a dyadic point is explicitly rational.
\end{lemma}
\begin{proof}
Conditioning on the independent tail gives the general, not just dyadic, identity
\[
 F(y)=\frac{2^{T_n}}{n!}
  \sum_{h=0}^{2^n-1}\eps_h
    \E\pos{y-2^{-n}(h+X)}^n.
\]
For $r\le n$, differentiate $r$ times. When $r=n$, the derivative of the truncated power is interpreted as $n!$ times the strict step function; the distribution of $X$ has no atoms, so its value at the cutoff is immaterial. At $y=m/2^n$, all terms $h<m$ have a nonnegative argument for every $X\in[0,1]$, whereas the remaining terms vanish. After replacing $X$ by $(1+Z)/2$, the scale is $2^{(n+1)r-T_n}$ and the expectation divided by $(n-r)!$ is $\Phi_{n-r}$. Higher derivatives vanish by \cref{lem:derivatives}.
\end{proof}

\subsection{An exact local deconvolution identity}
Let
\begin{equation}
 P_x(h)=\sum_{r=0}^{n}\frac{F^{(r)}(x)}{r!}h^r,
 \qquad x=\frac m{2^n}.
 \label{eq:dyadic-Taylor-poly}
\end{equation}
The polynomial $P_x$ is the Taylor polynomial at $x$, but no equality $F(x+h)=P_x(h)$ is asserted.

\begin{theorem}[Exact shrinking-cell formula]
\label{thm:local-cell}
Suppose $N\ge n\ge1$ and $x=m/2^n\in[0,1]$. For every real $h$ with $|h|\le2^{-N-1}$,
\begin{equation}
 \boxed{\displaystyle
 G_N(x+h)=
 \sum_{k=0}^{\lfloor n/2\rfloor}e_k\,2^{-2k(N+1)}
  \sum_{r=0}^{n-2k}\frac{F^{(r+2k)}(x)}{r!}h^r.}
 \label{eq:local-cell}
\end{equation}
All coefficients $e_k$ are the finite sums \cref{eq:e-partitions}. Equivalently, on this cell,
\begin{equation}
 G_N(x+h)=\bigl[M(2^{-N-1}D_h)\bigr]^{-1}P_x(h),
 \qquad D_h=\frac{\mathrm d}{\mathrm dh},
 \label{eq:local-operator}
\end{equation}
where the inverse operator means only its finite action on a polynomial of degree at most $n$.
\end{theorem}

\begin{proof}
Write $R=N-n$. The centered prefix splits into independent parts as
\[
 Y_N=S_n+2^{-n}Y_R,\qquad Y_R=\frac{1+Z_R}{2},
\]
where $Y_0=1/2$. The support of $Y_R$ is contained in
$[2^{-R-1},1-2^{-R-1}]$. The bound on $h$ implies
$|2^nh|\le2^{-R-1}$. Therefore the cutoff pattern in the first-$n$-variable box formula cannot change: all indices $\ell<m$ contribute full powers, and all indices $\ell\ge m$ contribute zero. At the two cell endpoints some arguments can equal zero, which causes no change because $n\ge1$. Thus
\begin{equation}
 G_N(x+h)=\frac{2^{-T_n}}{n!}
  \sum_{\ell=0}^{m-1}\eps_\ell
  \E(2m-2\ell-1+2^{n+1}h-Z_R)^n.
 \label{eq:local-tail-polynomial}
\end{equation}
Replacing $Z_R$ here by $Z$ produces a polynomial in $h$ whose $r$th derivative at zero is \cref{eq:derivative-master}. That polynomial is exactly $P_x(h)$.

For any polynomial $p(v)$, symmetry gives
$\E p(v-Z_R)=M_R(D_v)p(v)$ and
$\E p(v-Z)=M(D_v)p(v)$. The tail factorization says
\[
 M_R(t)=\frac{M(t)}{M(2^{-R}t)}.
\]
These identities are legitimate in the finite-dimensional space of polynomials: powers of $D_v$ eventually vanish, so no infinite convergence or analyticity of $F$ is being assumed. In \cref{eq:local-tail-polynomial}, $v=2m-2\ell-1+2^{n+1}h$, so $D_v=2^{-n-1}D_h$. The denominator is therefore
$M(2^{-R-n-1}D_h)=M(2^{-N-1}D_h)$.
Expanding its reciprocal by \cref{eq:inverse-M} gives \cref{eq:local-cell}.
\end{proof}

\begin{corollary}[Polynomial dependence on the squared tail scale]
\label{cor:scale-polynomial}
For a fixed dyadic point $x=m/2^n$, put $d=\lfloor n/2\rfloor$. Then, for all integers $N\ge n$,
\begin{equation}
 G_N(x)=Q_x(4^{-N}),\qquad
 Q_x(z)=\sum_{k=0}^{d}e_k4^{-k}F^{(2k)}(x)z^k.
 \label{eq:Qx}
\end{equation}
In particular $Q_x(0)=F(x)$ and $\deg Q_x\le d$.
\end{corollary}
This is an exact identity for the whole tail of the sequence $G_N(x)$, not a finite-order asymptotic expansion. Its constant term can therefore be extracted using finitely many depths.

\subsection{Why the degree bound is attained}
To show exact degree, it remains to rule out vanishing of the reciprocal coefficients.

\begin{lemma}[Signs of the reciprocal coefficients]
\label{lem:e-sign}
For every $k\ge0$, $(-1)^ke_k>0$.
\end{lemma}
\begin{proof}
Euler's sine product, in the standard normalization \cite[Eq.~4.22.1]{dlmf}, is
\[
 \frac{\sin t}{t}=\prod_{\ell=1}^{\infty}
    \left(1-\frac{t^2}{\pi^2\ell^2}\right).
\]
Near $t=0$, it gives
\begin{equation}
 \frac1{M(it)}=
 \prod_{j=1}^{\infty}\prod_{\ell=1}^{\infty}
 \left(1-\frac{t^2}{4^j\pi^2\ell^2}\right)^{-1}.
 \label{eq:e-positive-product}
\end{equation}
The sum of the positive numbers $(4^j\pi^2\ell^2)^{-1}$ is finite. Thus this product converges near zero, and its power-series coefficients are limits of the coefficients of finite products. Every coefficient of a finite product is nonnegative, and, for every degree, even one factor contributes a strictly positive coefficient. The coefficient of $t^{2k}$ on the left is $(-1)^ke_k$.
\end{proof}

\begin{proposition}[Exact degree at a reduced interior dyadic]
\label{prop:exact-degree}
If $0<m<2^n$, $m$ is odd, and $n\ge1$, then
\begin{equation}
 \deg Q_{m/2^n}=\floor{\frac n2}.
 \label{eq:exact-degree}
\end{equation}
\end{proposition}
\begin{proof}
When $n=2d$, the highest possible even derivative is
\[
 F^{(2d)}(m/2^{2d})=2^{T_{2d}}\Theta(m)
   =2^{T_{2d}}\eps_{(m-1)/2}\ne0.
\]
When $n=2d+1$, it is
\[
 F^{(2d)}(m/2^{2d+1})
 =2^{T_{2d}}\Theta(m/2)
 =2^{T_{2d}-1}\eps_{\lfloor m/4\rfloor}\ne0.
\]
This also covers $n=1$ and $d=0$. The coefficient of $z^d$ in \cref{eq:Qx} is nonzero by \cref{lem:e-sign}.
\end{proof}
A dyadic argument should therefore be reduced before applying the formula: the relevant depth is its actual denominator exponent, not a larger exponent used to represent the same number.

\section{Exact quarter-base extraction and its sharpness}
\label{sec:extraction}

\subsection{Arbitrary finite choices of depths}
Polynomial interpolation converts \cref{cor:scale-polynomial} directly into a formula free of Taylor data.

\begin{theorem}[Extraction from arbitrary distinct depths]
\label{thm:arbitrary-depths}
Fix $n\ge1$, $d=\lfloor n/2\rfloor$, and distinct integers $N_0,\ldots,N_d\ge n$. For every $x=m/2^n\in[0,1]$,
\begin{equation}
 F(x)=\sum_{i=0}^{d}W_iG_{N_i}(x),\qquad
 W_i=\prod_{\substack{0\le\ell\le d\\\ell\ne i}}
       \frac{4^{-N_\ell}}{4^{-N_\ell}-4^{-N_i}}.
 \label{eq:arbitrary-extraction}
\end{equation}
The weights depend only on the selected depths, not on $m$.
\end{theorem}
\begin{proof}
Let $z_i=4^{-N_i}$. For any polynomial $Q$ of degree at most $d$,
\[
 Q(z)=\sum_{i=0}^{d}Q(z_i)
       \prod_{\ell\ne i}\frac{z-z_\ell}{z_i-z_\ell}.
\]
To verify this identity without assuming interpolation theory, subtract the right-hand side. The difference has degree at most $d$ and vanishes at the $d+1$ distinct points $z_i$, so it is zero. Evaluating at $z=0$ and taking $Q=Q_x$ proves the theorem.
\end{proof}

\subsection{Consecutive depths and Gaussian binomial weights}
For completeness, our $q$-notation is
\begin{equation}
 (a;q)_r=\prod_{\ell=0}^{r-1}(1-aq^\ell),\qquad
 \qbinom{d}{j}{q}=\frac{(q;q)_d}{(q;q)_j(q;q)_{d-j}}.
 \label{eq:q-definitions}
\end{equation}
These are finite products. At $q=1/4$, all their values are rational.

Choose $N_j=n+R+j$. Their nodes are a common multiple of $q^j$, so that common multiple cancels from the weights. Splitting the factors in \cref{eq:arbitrary-extraction} into $\ell<j$ and $\ell>j$ gives
\begin{equation}
 \boxed{\displaystyle
 w_{d,j}(q)=
 \frac{(-1)^{d-j}q^{T_{d-j}}}{(q;q)_j(q;q)_{d-j}}
 =\frac{(-1)^{d-j}q^{T_{d-j}}}{(q;q)_d}
       \qbinom{d}{j}{q}.}
 \label{eq:q-weights}
\end{equation}
Indeed the lower-index factors are $1/(1-q^{j-\ell})$, and the upper-index factors are $-q^{\ell-j}/(1-q^{\ell-j})$. At $q=1/4$,
\[
 (q;q)_r=4^{-T_r}\prod_{s=1}^{r}(4^s-1),
\]
which turns \cref{eq:q-weights} into the integer-product formula \cref{eq:weights-integer}.

\begin{proof}[Completion of the proof of \cref{thm:main-quarter}]
Apply \cref{thm:arbitrary-depths} with $N_j=n+R+j$ and replace every $G_{N_j}(m/2^n)$ by its integer power sum \cref{eq:G-integer-sum}. The result is exactly \cref{eq:main-quarter-shift}; taking $R=0$ gives \cref{eq:main-quarter}.
\end{proof}

In particular, the weights satisfy the finite identities
\begin{equation}
 \sum_{j=0}^{d}w_{d,j}=1,\qquad
 \sum_{j=0}^{d}w_{d,j}4^{-jr}=0\quad(1\le r\le d).
 \label{eq:annihilation}
\end{equation}
These follow by extracting the value at zero of the polynomials $1,z,\ldots,z^d$. They express exactly which tail corrections are removed.

\subsection{A precise optimality statement}
The number $d+1$ should not be mistaken for a lower bound on all possible algorithms. Binary compression, for example, is not an interpolation algorithm. There is, however, a sharp lower bound within the natural class of uniform prefix-extraction rules.

\begin{theorem}[Minimal number of samples for a grid-wide linear rule]
\label{thm:minimal}
Fix $n\ge1$ and $d=\lfloor n/2\rfloor$. Suppose distinct depths $N_1,\ldots,N_L\ge n$ and real numbers $c_1,\ldots,c_L$, independent of $m$, satisfy
\[
 F(m/2^n)=\sum_{i=1}^{L}c_iG_{N_i}(m/2^n)
 \quad\text{for every }m=0,\ldots,2^n.
\]
Then $L\ge d+1$. Hence the consecutive-depth rule is optimal in this class.
\end{theorem}
\begin{proof}
The grid contains $x_0=1$ and $x_r=2^{-2r}$ for $1\le r\le d$. The polynomial $Q_{x_0}$ is the constant one. By \cref{prop:exact-degree}, $Q_{x_r}$ has degree exactly $r$. Therefore these $d+1$ polynomials form a basis of all polynomials of degree at most $d$: their leading terms make the coefficient matrix triangular with nonzero diagonal.

The supposed rule would consequently recover $Q(0)$ from $Q(z_i)$, $z_i=4^{-N_i}>0$, for every polynomial $Q$ of degree at most $d$. If $L\le d$, use $Q(z)=\prod_{i=1}^{L}(z-z_i)$. All sampled values are zero, but $Q(0)\ne0$, a contradiction. The case $L=0$ is included by the empty product $Q=1$.
\end{proof}
This theorem does not exclude weights depending on the point, nonlinear formulas, other kinds of samples, or a direct finite coefficient calculation.

\subsection{The extraction weights themselves are well conditioned}
Although the inclusion--exclusion sums can contain large cancellations, the final quarter-base interpolation row has a uniformly bounded absolute sum.

\begin{proposition}[Absolute-sum bound]
\label{prop:weight-bound}
For $0<q<1$,
\begin{equation}
 \sum_{j=0}^{d}|w_{d,j}(q)|
   =\frac{(-q;q)_d}{(q;q)_d}
   =\prod_{r=1}^{d}\frac{1+q^r}{1-q^r}.
 \label{eq:weight-norm}
\end{equation}
For $q=1/4$, this quantity is strictly less than $2$, uniformly in $d$.
\end{proposition}
\begin{proof}
The finite product identity
\[
 \prod_{r=1}^{d}(1+zq^r)
 =\sum_{\ell=0}^{d}z^\ell q^{T_\ell}\qbinom{d}{\ell}{q}
\]
can be verified by induction on $d$: multiply the identity for $d-1$ by $1+zq^d$ and put its two coefficient terms over the common finite-product denominator. At $z=1$, division by $(q;q)_d$ proves \cref{eq:weight-norm}.

For $q=1/4$, the first factor is $5/3$. For $r\ge2$, the elementary bound
$\log((1+u)/(1-u))\le2u/(1-u)$ gives
\[
 \sum_{r=2}^{\infty}\log\frac{1+4^{-r}}{1-4^{-r}}
 \le\frac{32}{15}\sum_{r=2}^{\infty}4^{-r}
 =\frac8{45}.
\]
Moreover $\log(6/5)\ge1/5-(1/5)^2/2=9/50>8/45$. Thus
$(5/3)\exp(8/45)<2$. For $d=0$ the absolute sum is one.
\end{proof}
If each accurately computed input $G_{N_j}(x)$ has absolute error at most $\delta$, combining them with the exact weights adds error less than $2\delta$. This does not bound the error incurred by computing each signed power sum in floating point; integer or rational arithmetic is preferable there.

\section{Eliminating moments by finitely many uniform prefixes}
\label{sec:moment-extraction}

The coefficients $a_k$ in the master and binary formulas admit a third type of recurrence-free representation. Instead of expanding the infinite product directly, interpolate a coefficient of its finite prefixes.

\subsection{Two explicit finite-prefix coefficients}
Write
\begin{equation}
 a_{k,N}=\coeff{t}{2k}M_N(t).
 \label{eq:prefix-coeff-definition}
\end{equation}
For $N=0$, set $a_{0,0}=1$ and $a_{k,0}=0$ for $k>0$. Expanding each of the $N$ factors gives the finite multinomial formula
\begin{equation}
 a_{k,N}=
 \sum_{\substack{r_1+\cdots+r_N=k\\r_i\ge0}}
 \prod_{i=1}^{N}\frac{4^{-ir_i}}{(2r_i+1)!}.
 \label{eq:prefix-coeff-multinomial}
\end{equation}
The sum is finite because every $r_i$ is at most $k$.

There is also a single signed power sum. For $N\ge1$ put $L_N=1-2^{-N}$. Factoring the hyperbolic sines yields
\begin{equation}
 M_N(t)=2^{N(N-1)/2}t^{-N}
 \sum_{h=0}^{2^N-1}\eps_h
        e^{(L_N-2^{1-N}h)t}.
 \label{eq:prefix-M-exponentials}
\end{equation}
Here the apparent pole at zero is removable: the numerator contains the product of $N$ factors each vanishing at zero. Expanding the exponentials gives
\begin{equation}
 \boxed{\displaystyle
 a_{k,N}=
 \frac{2^{N(N-1)/2-N(2k+N)}}{(2k+N)!}
 \sum_{h=0}^{2^N-1}\eps_h(2^N-1-2h)^{2k+N}.}
 \label{eq:prefix-coeff-power}
\end{equation}
To check the constant in \cref{eq:prefix-M-exponentials}, use
$H(t/2^j)=2^{j-1}t^{-1}(e^{t/2^j}-e^{-t/2^j})$ and then expand the finite product. The subset sums $\sum 2^{-j}$ run through $h/2^N$, with sign $\eps_h$; reversing the order of the $N$ binary digits does not change that sign.

\subsection{Exact coefficient extraction at the quarter base}
\begin{theorem}[Finite-prefix moment elimination]
\label{thm:moment-extraction}
For every $k\ge0$ and every integer $N_0\ge0$,
\begin{equation}
 \boxed{\displaystyle
 a_k=\sum_{j=0}^{k}w_{k,j}\,a_{k,N_0+j},}
 \label{eq:moment-extraction}
\end{equation}
where $w_{k,j}$ is the explicit rational weight \cref{eq:weights-integer}, and each $a_{k,N}$ on the right is given by \cref{eq:prefix-coeff-multinomial} or \cref{eq:prefix-coeff-power}. No infinite moments remain on the right.
\end{theorem}
\begin{proof}
From $M_N(t)=M(t)/M(2^{-N}t)$,
\begin{equation}
 a_{k,N}=\sum_{r=0}^{k}a_{k-r}e_r(4^{-N})^r.
 \label{eq:coefficient-scale-polynomial}
\end{equation}
For fixed $k$, this is a polynomial of degree at most $k$ in $4^{-N}$, whose constant term is $a_k$. Interpolation at the $k+1$ consecutive nodes proves the identity. The occurrence of $a_{k-r}$ in this proof is not an auxiliary recurrence in the final expression: all sampled coefficients on the right of \cref{eq:moment-extraction} have already been supplied by independent finite formulas.
\end{proof}
At $N_0=0$ and $k>0$, the $j=0$ term is zero. Thus only the nontrivial prefixes $1,\ldots,k$ are needed. For example,
\[
 a_{1,1}=\frac1{24},\qquad
 a_1=\frac43a_{1,1}=\frac1{18}.
\]
The raw moment $\E Z_1^2$ is $1/12$, not $a_{1,1}$; the normalization by $(2k)!$ remains essential.

Inserting \cref{eq:moment-extraction,eq:prefix-coeff-power} into \cref{eq:a-Phi} and then into either \cref{eq:master-Phi} or \cref{eq:binary-formula} produces another explicit family consisting entirely of finite Thue--Morse sums. It is distinct in organization from direct value extraction: the interpolation depth is controlled by the moment index $k$, and the recovered coefficient can be reused at many dyadic arguments.

\section{Rvachev values, the signed extension, and derivatives}
\label{sec:rvachev}

\subsection{A direct reduction for every dyadic Rvachev value}
For an integer $a$ and $n\ge0$, the support and folding identity give
\begin{equation}
 \boxed{\displaystyle
 \up\!\left(\frac a{2^n}\right)=
 \begin{cases}
 F\!\left(\dfrac{2^n-|a|}{2^n}\right),&|a|\le2^n,\\[2mm]
 0,&|a|>2^n.
 \end{cases}}
 \label{eq:up-dyadic-fold}
\end{equation}
Substituting any of the closed expressions above completes the evaluation, without a recurrence for $\up$ or $F$. At denominator depth zero this simply says $\up(0)=1$ and $\up(a)=0$ at every other integer.

One can also extract the density directly from finite box splines. Let $U_N$ denote the density of $Z_N$, with $N\ge2$ so it is continuous. Differentiating the finite box CDF gives
\begin{equation}
 U_N(x)=\frac{2^{N(N-1)/2}}{(N-1)!}
   \sum_{h=0}^{2^N-1}\eps_h
   \pos{x+1-2^{-N}-2^{1-N}h}^{N-1}.
 \label{eq:up-finite-spline}
\end{equation}
Every summand is an explicit rational number at rational $x$.

\begin{proposition}[Direct finite-spline formula for $\up$]
\label{prop:up-extraction}
For $n\ge1$, $|a|\le2^n$, $d=\lfloor n/2\rfloor$, and $R\ge0$,
\begin{equation}
 \boxed{\displaystyle
 \up\!\left(\frac a{2^n}\right)
 =\sum_{j=0}^{d}w_{d,j}\,
 U_{n+1+R+j}\!\left(\frac a{2^n}\right).}
 \label{eq:up-extraction}
\end{equation}
Together with \cref{eq:up-finite-spline}, this is a finite formula for the Rvachev value itself.
\end{proposition}
\begin{proof}
The independent-copy identity $Z_N\overset d=(V+Z_{N-1}')/2$ gives
\[
 U_N(x)=\Pp(2x-1\le Z_{N-1}'\le2x+1).
\]
The finite law is symmetric and supported in $[-1,1]$. For $x\ge0$, its upper cutoff can be discarded; symmetry then gives
$U_N(x)=\Pp(Z_{N-1}'\le1-2x)=G_{N-1}(1-x)$.
Evenness yields the identity
\begin{equation}
 U_N(x)=G_{N-1}(1-|x|)\quad(x\in\R).
 \label{eq:finite-fold}
\end{equation}
Apply the quarter-base formula at $(2^n-|a|)/2^n$, using depths $n+R+j$, and fold back.
\end{proof}
The shift by one in \cref{eq:up-extraction} is important: $U_N$ is a density of a sum of $N$ uniforms, whereas $G_{N-1}$ is a CDF for $N-1$ uniforms.

\subsection{The signed global extension without functional recursion}
For $m\ge0$, define
\begin{equation}
 b=\floor{\frac{m}{2^{n+1}}},\qquad
 r=m-2^{n+1}b,\qquad
 a=\min\{r,2^{n+1}-r\}.
 \label{eq:global-fold-data}
\end{equation}
Then $0\le a\le2^n$, and the definition of $\Theta$ gives
\begin{equation}
 \boxed{\displaystyle
 \Theta\!\left(\frac m{2^n}\right)
 =(-1)^{\bits(b)}F\!\left(\frac a{2^n}\right).}
 \label{eq:global-dyadic}
\end{equation}
For $m<0$ the value is zero. Indeed, on the interval $[2b,2b+2]$, one has
$\Theta(x)=\eps_b\up(x-2b-1)$, and folding the bump around its center gives precisely $a/2^n$. Thus signed values at arbitrarily large positive dyadic arguments require only binary digits and one principal-interval closed formula.

The globally bounded distribution function instead has the constant values $F(x)=0$ for $x\le0$ and $F(x)=1$ for $x\ge1$. These two extensions should not be interchanged.

\subsection{A companion closed form for derivatives}
The differentiated master formula \cref{eq:derivative-master} already evaluates $F^{(r)}$ at every dyadic point. For $x\in[-1,1]$, the identity $\up(x)=\Theta(x+1)$ gives
\begin{equation}
 \up^{(r)}(x)=2^{T_r}\Theta\bigl(2^r(x+1)\bigr).
 \label{eq:up-derivatives}
\end{equation}
At dyadic $x$, the right-hand side is evaluated by \cref{eq:global-dyadic} and a closed formula for $F$. For $r\ge n+1$ and $x=a/2^n$, it is zero. Outside $[-1,1]$ all derivatives vanish. At the boundary points the same conclusion follows from smoothness and flatness. Hence the same finite arithmetic also supplies all finite Taylor jets needed in the local spline theorem.

\section{An elementary common denominator and a rounding formula}
\label{sec:denominator}

The next formula uses an integer-part operation in addition to finite arithmetic. It is included as a different exact mechanism, not as a replacement for the purely algebraic formulas above.

\begin{theorem}[An explicit common denominator]
\label{thm:denominator}
For $n\ge0$ set $d=\lfloor n/2\rfloor$ and
\begin{equation}
 \mathcal D_n=2^{T_n}(n+d)!\prod_{j=1}^{d}(4^j-1).
 \label{eq:common-denominator}
\end{equation}
Then
\begin{equation}
 \mathcal D_nF(m/2^n)\in\mathbb Z
 \qquad(0\le m\le2^n).
 \label{eq:denominator-integrality}
\end{equation}
\end{theorem}
\begin{proof}
Use the master formula with the positive composition expression \cref{eq:Phi-compositions}. A typical summand corresponding to $k\le d$ and a composition $k=r_1+\cdots+r_\ell$ has a denominator dividing
\[
 2^{T_n}(n-2k)!
  \prod_{i=1}^{\ell}(2r_i+1)!
  \prod_{i=1}^{\ell}(4^{R_i}-1),
 \qquad R_i=r_i+\cdots+r_\ell.
\]
The suffix sums $R_i$ are distinct integers between $1$ and $k$, so the last product divides $\prod_{j=1}^{d}(4^j-1)$. The sum of the factorial arguments is
\[
 (n-2k)+\sum_{i=1}^{\ell}(2r_i+1)=n+\ell\le n+d.
\]
By the integrality of multinomial coefficients, their factorial product divides $(n+\ell)!$, which divides $(n+d)!$. The term $k=0$ has just the factorial $n!$ and obeys the same bound. Every summand, and therefore their integer-signed sum, has denominator dividing $\mathcal D_n$.
\end{proof}
This bound is intentionally elementary and is not asserted to be a least common denominator or to improve the sharper arithmetic bounds in \cite{arias2018}. Its merit here is that it follows immediately from a recurrence-free positive expression.

\begin{lemma}[Uniform prefix error]
\label{lem:prefix-error}
For every $N\ge1$ and every real $x$,
\begin{equation}
 |G_N(x)-F(x)|\le2^{-N}.
 \label{eq:prefix-error}
\end{equation}
\end{lemma}
\begin{proof}
Couple the variables by
$X=Y_N+2^{-N-1}Z'$, where $Z'$ is an independent copy of $Z$ and $|Z'|\le1$. Writing $\delta=2^{-N-1}$, this gives
\[
 F(x-\delta)\le G_N(x)\le F(x+\delta).
\]
The $2$-Lipschitz bound for $F$ proves the claim.
\end{proof}

\begin{corollary}[Certified one-prefix rounding]
\label{cor:rounding}
For $n\ge1$, choose the explicit integer
\begin{equation}
 N=\max\left\{n,\ 2+\floor{\log_2\mathcal D_n}\right\}.
 \label{eq:round-depth}
\end{equation}
Then
\begin{equation}
 \boxed{\displaystyle
 F\!\left(\frac m{2^n}\right)=
 \frac1{\mathcal D_n}
 \floor{\frac{\mathcal D_n\mathcal S_N(2^{N-n}m)}
             {2^{T_N}N!}+\frac12}.}
 \label{eq:rounding-formula}
\end{equation}
\end{corollary}
\begin{proof}
The selected depth satisfies $2^N>2\mathcal D_n$. Therefore
$|\mathcal D_nG_N(x)-\mathcal D_nF(x)|<1/2$, and the latter scaled value is an integer by \cref{thm:denominator}. Nearest-integer rounding is exact, with no tie. Substitute \cref{eq:G-integer-sum}.
\end{proof}
The base-two logarithm in \cref{eq:round-depth} can be replaced by the binary digit length of the explicitly specified positive integer $\mathcal D_n$. This avoids any floating-point logarithm. Although the formula has only one finite prefix, that prefix is usually enormous. It is a mathematical certification device, not a recommended large-depth implementation.

\section{Worked examples and practical evaluation}
\label{sec:examples}

\subsection{Two samples recover the quarter value}
At $x=1/4$, take $n=2$ and $d=1$. The explicit spline sums give
\[
 G_2(1/4)=\frac1{16},\qquad G_3(1/4)=\frac{13}{192}.
\]
For example, the numerator for $G_3(1/4)$ is $3^3-1^3=26$ and its denominator is $2^{T_3}3!=384$. Therefore
\begin{equation}
 F(1/4)=\frac{-G_2(1/4)+4G_3(1/4)}3
       =\frac5{72}.
 \label{eq:example-quarter}
\end{equation}
The exact scale polynomial is
\begin{equation}
 G_N(1/4)=\frac5{72}-\frac19\,4^{-N}\qquad(N\ge2).
 \label{eq:quarter-poly}
\end{equation}
Indeed $F'(1/4)=1$, $F''(1/4)=8$, and $e_1=-1/18$. The full local statement is
\begin{equation}
 G_N(1/4+h)=\frac5{72}+h+4h^2-\frac19\,4^{-N},
 \qquad |h|\le2^{-N-1}.
 \label{eq:quarter-cell}
\end{equation}
After shifting the prefix index by one, this is the quarter-point cell recorded in the repository \cite{repo}. The general shrinking-cell theorem explains the same phenomenon at every dyadic point and at all orders.

\subsection{Three samples recover a sixteenth value}
At $x=1/16$, use $n=4$ and $d=2$. Direct integer sums give
\[
 G_4(1/16)=\frac1{24576},\qquad
 G_5(1/16)=\frac{121}{1966080},\qquad
 G_6(1/16)=\frac{6331}{94371840}.
\]
The exact value is
\begin{equation}
 F(1/16)=\frac{G_4(1/16)-20G_5(1/16)+64G_6(1/16)}{45}
        =\frac{143}{2073600}.
 \label{eq:example-sixteenth}
\end{equation}
Here the polynomial in $t=4^{-N}$ is genuinely quadratic:
\begin{align}
 G_N(1/16)&=\frac{143}{2073600}-\frac5{648}t
                         +\frac{248}{2025}t^2,\label{eq:sixteenth-poly}\\
 G_N(3/16)&=\frac{46657}{2073600}-\frac{67}{648}t
                         -\frac{248}{2025}t^2,
 \qquad N\ge4.\label{eq:three-sixteenth-poly}
\end{align}
For comparison, the eighth value still requires only two samples:
\[
 G_3(1/8)=\frac1{384},\qquad G_4(1/8)=\frac5{1536},
 \qquad F(1/8)=\frac1{288}.
\]
Its exact scale polynomial is $G_N(1/8)=1/288-4^{-N}/18$.

\subsection{A numerator with three nonzero bits}
Consider $x=21/64$. The binary numerator is $21=(10101)_2$, so $J(21)=\{0,2,4\}$. The block parameters give
\begin{equation}
 F(21/64)=2^{-3}\Phi_2(13/8)
          -2^{-10}\Phi_4(3/2)
          +2^{-21}\Phi_6(1).
 \label{eq:binary-example}
\end{equation}
The three terms, evaluated with the positive composition coefficients, are respectively
\[
 \frac{1585}{9216},\qquad
 -\frac{71179}{265420800},\qquad
 \frac{1153}{561842749440}.
\]
Their sum is
\begin{equation}
 F(21/64)=\frac{482385079229}{2809213747200}.
 \label{eq:value-21-64}
\end{equation}
This illustrates the distinction between binary compression and interpolation: the former reduces the number of blocks according to the numerator, whereas the latter uses a uniform row determined by the denominator depth.

\subsection{A table of exact values}
All entries in \cref{tab:values} were independently matched by the master, binary, and quarter-base formulas using rational arithmetic. Reflection supplies the remaining values on each displayed grid. For example, $\up(3/4)=F(1/4)=5/72$ and $\up(5/8)=F(3/8)=73/288$.

\begin{table}[htbp]
\centering
\caption{Selected exact Fabius values. Fractions are reduced.}
\label{tab:values}
\renewcommand{\arraystretch}{1.3}
\begin{tabular}{@{}rl@{\qquad}rl@{}}
\toprule
$x$ & $F(x)$ & $x$ & $F(x)$\\
\midrule
$1/2$ & $1/2$ & $1/4$ & $5/72$\\
$1/8$ & $1/288$ & $3/8$ & $73/288$\\
$1/16$ & $143/2073600$ & $3/16$ & $46657/2073600$\\
$5/16$ & $305857/2073600$ & $7/16$ & $777743/2073600$\\
$1/32$ & $19/33177600$ & $13/32$ & $10393219/33177600$\\
\bottomrule
\end{tabular}
\end{table}
Two further examples with larger denominators are
\[
 F(21/64)=\frac{482385079229}{2809213747200},\qquad
 F(37/128)=\frac{4121049919811}{35957935964160}.
\]

\subsection{Operation counts and exact arithmetic}
For the consecutive-depth quarter formula with $R=0$, the raw integer sums contain in total
\begin{equation}
 \sum_{j=0}^{d}2^jm=m(2^{d+1}-1)
 \label{eq:power-term-count}
\end{equation}
integer-power summands. This is a count before exploiting symmetry, cancellation identities, or reuse. The largest depth is $n+\lfloor n/2\rfloor$. Reducing the dyadic fraction first can substantially decrease this count. Replacing $x$ by $1-x$ when it is closer to zero also decreases the numerator, using reflection at the end.

For the binary formula, only $\bits(m)$ outer blocks are present. The required coefficient indices are at most $\lfloor n/2\rfloor$, regardless of the numerator. There are $2^{k-1}$ positive ordered compositions of $k\ge1$: choose cuts among $k-1$ possible positions. The multiplicity formula instead has one term per integer partition of $k$. Coefficients computed from either finite expression may be cached without changing their recurrence-free definition.

The signs in $\mathcal S_N(A)$ can produce severe cancellation between very large integers. Compute that numerator exactly, divide only after summation, and combine interpolation weights as exact fractions. The bounded norm of the weight row does not make a low-precision evaluation of the individual sums safe. Conversely, a carefully designed recursive algorithm may outperform all fully expanded sums. No claim that a closed form must be the fastest algorithm is intended.

\section{A related integer-base extension}
\label{sec:base-extension}

The squared-scale mechanism is not restricted to base two. The extension in this section concerns a \emph{different family of distributions}; it is not a claim about rationality of the original Fabius function at nondyadic rational arguments.

For an integer $b\ge2$, let
\begin{equation}
 X_b=(b-1)\sum_{j=1}^{\infty}b^{-j}U_j,
 \qquad F_b(x)=\Pp(X_b\le x).
 \label{eq:base-X}
\end{equation}
The normalization keeps the support equal to $[0,1]$. Write
\begin{equation}
 \begin{split}
 Z_b&=2X_b-1=(b-1)\sum_{j\ge1}b^{-j}V_j,\\
 Y_{b,N}&=(b-1)\sum_{j=1}^{N}b^{-j}U_j+\frac1{2b^N},
 \qquad G_{b,N}(x)=\Pp(Y_{b,N}\le x).
 \end{split}
 \label{eq:base-prefix}
\end{equation}
The same characteristic-function argument as before gives a smooth density for the infinite sum. Its centered moment generating function is
\[
 M_b(t)=\prod_{j\ge1}H\!\left(\frac{(b-1)t}{b^j}\right).
\]
It is even, and its logarithmic coefficients have the finite description
\begin{equation}
 \log M_b(t)=\sum_{r\ge1}\lambda_{r,b}t^{2r},\qquad
 \lambda_{r,b}=
 \frac{(b-1)^{2r}2^{2r-1}B_{2r}}
 {r(2r)!(b^{2r}-1)}.
 \label{eq:base-lambda}
\end{equation}
The proof is the same geometric summation of $\log H$ used in \cref{sec:cumulants}; \cref{eq:Bernoulli-finite} supplies the Bernoulli numbers explicitly.

For finite $N\ge1$, inclusion--exclusion gives the rational formula
\begin{equation}
\begin{split}
 G_{b,N}(x)={}&\frac{b^{T_N}}{(b-1)^NN!}
  \sum_{(\delta_1,\ldots,\delta_N)\in\{0,1\}^N}
  (-1)^{\delta_1+\cdots+\delta_N}\\[-1mm]
 &\hspace{8mm}\cdot
 \pos{x-\frac1{2b^N}-(b-1)\sum_{j=1}^{N}b^{-j}\delta_j}^{N}.
\end{split}
\label{eq:base-spline}
\end{equation}
All sums and products are explicitly bounded.

\begin{theorem}[Exact extraction at $q=b^{-2}$]
\label{thm:base-extraction}
Let $n\ge1$, $0\le m\le b^n$, $d=\lfloor n/2\rfloor$, and $R\ge0$. Then
\begin{equation}
 \boxed{\displaystyle
 F_b(m/b^n)=\sum_{j=0}^{d}w_{d,j}(b^{-2})
               G_{b,n+R+j}(m/b^n),}
 \label{eq:base-extraction}
\end{equation}
where the weights are the finite products in \cref{eq:q-weights}. Consequently every value $F_b(m/b^n)$ is rational.
\end{theorem}
\begin{proof}
Split the first $n$ uniforms from the tail. After multiplication by $b^n$, every knot contributed by the first $n$ uniforms is an integer of the form
$(b-1)\sum_{j=1}^{n}b^{n-j}\delta_j$. The infinite rescaled tail is supported in $[0,1]$. Thus, at an integer grid point $m$, each truncated-power term is either a full power or zero. For the centered prefix at depth $N\ge n$, the same cutoff pattern remains valid throughout $|h|\le1/(2b^N)$, because the rescaled finite tail has its support inset by exactly that half-scale after restoring the original units.

Repeating the polynomial argument of \cref{thm:local-cell}, let $P_{b,x}(h)$ be the degree-$n$ polynomial obtained with the infinite tail. The finite-tail moment generating function is $M_b(t)/M_b(b^{-(N-n)}t)$. After converting to derivatives in the original variable $h$, it follows that
\[
 G_{b,N}(x+h)=
 \bigl[M_b(\tfrac12b^{-N}D_h)\bigr]^{-1}P_{b,x}(h)
 \quad\left(|h|\le\frac1{2b^N}\right).
\]
As before, this is a finite polynomial-operator identity, not an infinite operation on $F_b$. The reciprocal of $M_b$ is even. Setting $h=0$ therefore shows that $G_{b,N}(x)$ is a polynomial of degree at most $d$ in $b^{-2N}$ with constant term $P_{b,x}(0)=F_b(x)$. Interpolation at the consecutive nodes proves \cref{eq:base-extraction}. Substituting \cref{eq:base-spline} shows rationality directly.
\end{proof}
For example, $F_3(1/9)=7/576$, while $F_2=F$. The existence of exact finite extraction is caused by two structures together: grid-aligned finite-convolution knots and an even centered tail. Either structure alone would not imply the displayed degree bound.

\section{What the alternatives establish}
\label{sec:conclusion}

The closed-form requirement can be met in several genuinely different ways. Direct quarter-base extraction removes all moments and derivatives at once. Binary compression reduces the outer arithmetic to the numerator's nonzero digits, while positive composition or Bernoulli multiplicity sums remove its auxiliary coefficients. Finite-prefix coefficient extraction gives a third way to eliminate those coefficients. A common-denominator bound together with a rigorous prefix error yields an additional integer-part formula.

The central structural result is stronger than rationality: a centered finite spline, at a fixed dyadic point, lies on an exact polynomial in its squared tail scale. Its shrinking-cell version identifies all local polynomial coefficients simultaneously. This explains why a finite interpolation row can recover the limiting value exactly, why the quarter base replaces the half base, and why the necessary grid-wide number of samples is $\lfloor n/2\rfloor+1$.

All displayed value formulas terminate. Infinite products and probability models appear in the proofs, but none need to be approximated when evaluating a specified dyadic value. The article does not require an unproved conjecture, does not assert rationality at arbitrary rational arguments, and does not conflate a terminating dyadic Taylor jet with local analyticity of the Fabius function.

\clearpage
\appendix
\section{A compact implementation of the principal formula}
\label{app:code}

The following self-contained Python implementation uses only integer powers, factorials, binary digit counts, and exact fractions. No values of $F$, moments, or interpolation coefficients are generated by a recurrence. The function accepts unreduced dyadic representations as well as reduced ones. Python 3.10 or later is sufficient.

\begin{lstlisting}[language=Python]
from fractions import Fraction
from math import factorial, prod


def fabius_dyadic(m: int, n: int) -> Fraction:
    """Exact bounded F(m / 2**n), for 0 <= m <= 2**n."""
    if not isinstance(m, int) or not isinstance(n, int):
        raise TypeError("m and n must be integers")
    if n < 0 or not 0 <= m <= 2**n:
        raise ValueError("Require n >= 0 and 0 <= m <= 2**n")
    if n == 0:
        return Fraction(m)
    d = n // 2
    answer = Fraction(0)
    for j in range(d + 1):
        numerator = (-1)**(d-j) * 4**(j*(j+1)//2)
        denominator = (
            prod(4**r - 1 for r in range(1, j+1))
            * prod(4**r - 1 for r in range(1, d-j+1))
        )
        weight = Fraction(numerator, denominator)
        N = n + j
        A = m * 2**j
        power_sum = sum(
            (-1)**h.bit_count() * (2*A - 2*h - 1)**N
            for h in range(A)
        )
        prefix = Fraction(
            power_sum, 2**(N*(N+1)//2) * factorial(N)
        )
        answer += weight * prefix
    return answer


def rvachev_dyadic(a: int, n: int) -> Fraction:
    """Exact up(a / 2**n) on the whole real dyadic grid."""
    if not isinstance(a, int) or not isinstance(n, int):
        raise TypeError("a and n must be integers")
    if n < 0:
        raise ValueError("n must be nonnegative")
    if abs(a) > 2**n:
        return Fraction(0)
    return fabius_dyadic(2**n - abs(a), n)


assert fabius_dyadic(1, 4) == Fraction(143, 2073600)
assert rvachev_dyadic(3, 2) == Fraction(5, 72)
\end{lstlisting}

For larger depths, the companion file \texttt{verify\_closed\_forms.py} also contains the binary formula, composition and multiplicity formulas, finite-prefix moment extraction, the local-cell evaluator, and direct Rvachev spline extraction. Caching those explicit coefficients can improve repeated evaluations without changing any mathematical definition.

\section{Exact verification and reproducibility}
\label{app:verification}

The accompanying verification script uses \texttt{fractions.Fraction} and Python integers throughout. Its recorded run has status \texttt{PASS}. The tests are summarized in \cref{tab:checks}; the machine-readable record is \texttt{verification\_results.json}.

\begin{table}[htbp]
\centering
\caption{Exact checks performed by the accompanying script.}
\label{tab:checks}
\begin{tabularx}{\textwidth}{@{}Xr@{}}
\toprule
Check & Cases\\
\midrule
Master, binary, and quarter-base formulas on all grid representations $m/2^n$, $0\le n\le8$ & 520\\
Shifted-depth extractions, using $R=1,2$ at selected points & 88\\
Local-cell identity at five offsets, on grids through depth five and three prefix depths & 1005\\
Direct Rvachev spline extraction and the finite folding identity & 134\\
Moment coefficients $a_k$, $0\le k\le8$, by composition, multiplicity, prefix extraction, and a separate oracle & 9\\
Common-denominator integrality on the full tested grids & 520\\
Certified rounding on grids of depths one through three & 17\\
Integer-base formulas for $b=2,3,4$, including shifted-depth consistency and reflection & 146\\
\bottomrule
\end{tabularx}
\end{table}

The 520 cases are grid \emph{representations}, not 520 distinct rational points: for example, $1/2$ is tested again as $2/4$. The script also checks reflection on those grids, the exact-degree nonvanishing at every tested reduced interior dyadic, and the interpolation annihilation identities for $0\le d\le8$.

An independent moment oracle is deliberately isolated in the test code. It uses the standard triangular relation obtained from $M(2t)=H(t)M(t)$ solely to cross-check the finite formulas. No closed-form evaluator calls this oracle, and none requires recursively generated moments. For bases three and four, the reported checks are consistency checks between two extraction depths and reflection, not an independent characterization of those probability laws. The proofs supply that characterization.

To reproduce the recorded tests, run
\begin{lstlisting}[language=bash]
python verify_closed_forms.py --max-depth 8 \
  --output verification_results.json
\end{lstlisting}
The exhaustive signed sums grow rapidly with the requested maximum depth. The tests support the algebra and catch indexing and normalization mistakes; they are not substitutes for the all-depth proofs in the article.

The LaTeX source has no external figure or bibliography dependency. Compile it twice with \texttt{pdflatex}; a third pass may be used if a local installation still requests updated cross-references. The source uses Libertinus when available and otherwise falls back to Latin Modern. The archive does not contain font files.

\section{Source scope and relation to prior formulas}
\label{app:sources}

The principal repository document was consulted on its live \texttt{main} branch on 4 September 2026. That is a source-access description, not a commit-pinned or repository-wide audit. It provides the normalization, established dyadic arithmetic background, and several finite-coefficient and $q$-series precedents. The article is therefore organized around a different exact-extraction mechanism rather than presenting every moment identity as a new discovery.

The arithmetic treatment of Arias de Reyna supplies the classical rationality and denominator context. Haugland treats dyadic evaluation, while the explicit-formula questions of Reshetnikov and Pinelis show the relevance of finite products and uniform convolutions. In particular, a finite-prefix expression followed by an infinite limit does not by itself meet the closed-form convention used here. The quarter-base identity replaces that limit by a fixed finite interpolation row whose size is specified by the denominator depth.

The strengthened local-cell statement, exact-degree argument, uniform-rule minimality statement, and binary block formula are proved in full above. They are presented as the results of this derivation, without claiming that no equivalent statement exists elsewhere. Nothing in the accompanying files is represented as a completed Lean formalization.

\begin{thebibliography}{9}

\bibitem{repo}
V. Reshetnikov and contributors,
\emph{Fabius Function and Rvachev Up},
ProveIt repository, \path{Analysis/FabiusFunction/docs/Fabius_Function_and_Rvachev_Up},
live \texttt{main} branch consulted 4 September 2026.
\url{https://github.com/VladimirReshetnikov/ProveIt/tree/main/Analysis/FabiusFunction/docs/Fabius_Function_and_Rvachev_Up}.

\bibitem{arias2018}
J. Arias de Reyna,
\emph{Arithmetic of the Fabius function},
Integers \textbf{18} (2018), Article A51.
\url{https://math.colgate.edu/~integers/s51/s51.pdf}.
Preprint: \url{https://arxiv.org/abs/1702.06487}.

\bibitem{arias1982}
J. Arias de Reyna,
\emph{An infinitely differentiable function with compact support: Definition and properties},
English translation of the 1982 paper, arXiv:1702.05442 (2017).
\url{https://arxiv.org/abs/1702.05442}.

\bibitem{haugland}
J. K. Haugland,
\emph{Evaluating the Fabius function},
arXiv:1609.07999, version 2 (2020; first version 2016).
\url{https://arxiv.org/abs/1609.07999v2}.

\bibitem{reshetnikov}
V. Reshetnikov,
\emph{Conjectured formula for the Fabius function},
Mathematics Stack Exchange, question posted 5 July 2019.
\url{https://math.stackexchange.com/questions/3283519/conjectured-formula-for-the-fabius-function}.

\bibitem{pinelis}
I. Pinelis,
Answer to \emph{A non-recursive, explicit formula for the Fabius function},
MathOverflow, 23 January 2020.
\url{https://mathoverflow.net/questions/351004/a-non-recursive-explicit-formula-for-the-fabius-function}.

\bibitem{dlmf}
NIST Digital Library of Mathematical Functions,
\S4.22, Infinite Products and Partial Fractions, equation 4.22.1.
\url{https://dlmf.nist.gov/4.22.E1}.

\end{thebibliography}
\end{document}
